\chapter{พื้นฐานสถาปัตยกรรม OpenXR และระบบประมวลผลเชิงพื้นที่สำหรับสะเต็มศึกษา}
\label{chap:ch01}
\index{OpenXR}
\index{สะเต็มศึกษา (STEM Education)}
\index{สถาปัตยกรรม OpenXR}

% ==========================================================
% แผนบริหารการสอนประจำบทที่ 1
% ==========================================================
\section*{แผนบริหารการสอนประจำบทที่ 1}
\addcontentsline{toc}{section}{แผนบริหารการสอนประจำบทที่ 1}

\noindent\textbf{หัวข้อเนื้อหาประจำบท}
\begin{enumerate}
    \item วิวัฒนาการของเทคโนโลยีความเป็นจริงขยายและความท้าทายในสะเต็มศึกษา
    \item โครงสร้างสถาปัตยกรรมมาตรฐานเปิด OpenXR และลำดับชั้นการทำงาน
    \item การจัดการวงรอบชีวิตของระบบ (Instance, System, Session, and Frame Loop)
    \item การประยุกต์ใช้การประมวลผลเชิงพื้นที่ในการสร้างห้องทดลองวิทยาศาสตร์เสมือนจริง
\end{enumerate}

\noindent\textbf{วัตถุประสงค์เชิงพฤติกรรม}
\begin{enumerate}
    \item อธิบายปัญหาความแตกแยกของระบบและบทบาทของมาตรฐาน OpenXR ได้ถูกต้อง
    \item จำแนกหน้าที่ของ Application Layer, Loader, และ Runtime ได้อย่างแม่นยำ
    \item วิเคราะห์สถานะและการเปลี่ยนสถานะในวงรอบ Session และ Frame Loop ได้
    \item สรุปประโยชน์ของ OpenXR ในการพัฒนาสื่อการเรียนรู้สะเต็มศึกษาข้ามแพลตฟอร์มได้
\end{enumerate}

\noindent\textbf{กิจกรรมการเรียนการสอน}
\begin{enumerate}
    \item การบรรยายเชิงปฏิสัมพันธ์พร้อมการสาธิตการทำงานของ OpenXR Architecture
    \item กิจกรรมการวิเคราะห์เปรียบเทียบสถาปัตยกรรม Proprietary SDK กับ OpenXR Standard
    \item การฝึกปฏิบัติติดตั้งและกำหนดค่าสภาพแวดล้อมจำลอง Monado และ OpenXR Loader
    \item การอภิปรายกลุ่มย่อยเกี่ยวกับการนำเทคโนโลยีเชิงพื้นที่ไปใช้แก้ปัญหาการทดลองจริง
\end{enumerate}

\noindent\textbf{สื่อการเรียนการสอน}
\begin{enumerate}
    \item เอกสารคำสอนและสไลด์บรรยายอิเล็กทรอนิกส์ประจำบทที่ 1
    \item แผนผังสถาปัตยกรรม OpenXR Specification จาก Khronos Group
    \item ชุดทดสอบ Monado OpenXR Runtime และแอปพลิเคชันตัวอย่าง Hello XR
    \item ระบบการจัดการเรียนรู้ดิจิทัล (LMS) สำหรับการส่งงานและอภิปรายผล
\end{enumerate}

\noindent\textbf{การประเมินผล}
\begin{enumerate}
    \item การประเมินความรู้ผ่านแบบทดสอบย่อยเชิงมโนทัศน์ประจำบทที่ 1
    \item การประเมินผลการฝึกปฏิบัติการติดตั้งและทดสอบ OpenXR Runtime
    \item การประเมินการมีส่วนร่วมในการอภิปรายและการตอบคำถามในชั้นเรียน
\end{enumerate}

\newpage

\section{วิวัฒนาการและความท้าทายในสะเต็มศึกษา}

การจัดการศึกษาวิทยาศาสตร์และเทคโนโลยีในปัจจุบันเผชิญกับข้อจำกัดหลายประการ การทดลองทางวิทยาศาสตร์จำนวนมากเกี่ยวข้องกับปรากฏการณ์ที่มีอันตรายสูง สารเคมีที่มีพิษ อุปกรณ์ฟิสิกส์พลังงานสูงที่มีราคาแพง หรือปรากฏการณ์ระดับจุลภาคและดาราศาสตร์ที่ไม่สามารถสังเกตได้ด้วยตาเปล่า การนำเทคโนโลยีความเป็นจริงขยาย (Extended Reality หรือ XR) มาประยุกต์ใช้ จึงเป็นหนึ่งในกลยุทธ์สำคัญที่เข้ามาเติมเต็มกระบวนการเรียนรู้เชิงรุก (Active Learning)

อย่างไรก็ดี ในช่วงทศวรรษที่ผ่านมา นักพัฒนาซอฟต์แวร์การศึกษาต้องเผชิญกับอุปสรรคสำคัญจากปัญหาความแตกแยกของระบบ (Ecosystem Fragmentation) ดังแสดงในตารางที่ 1.1

\begin{table}[htbp]
\centering
\caption{การเปรียบเทียบระหว่างระบบดั้งเดิมเฉพาะของผู้ผลิตกับมาตรฐานเปิด OpenXR}
\label{tab:openxr_comparison}
\begin{tabularx}{\textwidth}{lYY}
\toprule
\textbf{มิติการเปรียบเทียบ} & \textbf{ระบบดั้งเดิมเฉพาะราย (Proprietary SDKs)} & \textbf{มาตรฐานเปิด (Khronos OpenXR)} \\
\midrule
การพกพาของโค้ด & ต้องพัฒนาโค้ดใหม่แยกตามยี่ห้อฮาร์ดแวร์ & พัฒนาครั้งเดียวสามารถรันได้ทุกอุปกรณ์ \\
สถาปัตยกรรมไดรเวอร์ & เชื่อมต่อกับฮาร์ดแวร์โดยตรงผ่านไดรเวอร์เฉพาะ & ผ่านเลเยอร์ OpenXR Loader และ Runtime \\
การบำรุงรักษา & มีภาระในการอัปเดต SDK หลายตัวขนานกัน & มีมาตรฐาน API กลางที่สม่ำเสมอและมั่นคง \\
การพัฒนาสื่อสะเต็ม & จำกัดเฉพาะอุปกรณ์ที่สถานศึกษามีงบจัดซื้อ & ใช้งานได้กับทั้งแว่น VR/AR มือถือ และคอมพิวเตอร์ \\
\bottomrule
\end{tabularx}
\end{table}

\section{โครงสร้างสถาปัตยกรรมมาตรฐานเปิด OpenXR}

มาตรฐาน OpenXR ซึ่งได้รับการพัฒนาและกำกับดูแลโดย Khronos Group ถูกออกแบบมาเพื่อแก้ปัญหาดังกล่าว โดยแบ่งโครงสร้างออกเป็น 3 ส่วนหลัก ได้แก่ Application Layer, OpenXR Loader, และ OpenXR Runtime

\begin{figure}[htbp]
\centering
\begin{tikzpicture}[node distance=1.5cm, auto]
    % Style definitions
    \tikzstyle{appbox} = [rectangle, rounded corners, minimum width=9cm, minimum height=1cm, text centered, draw=darkNavy, fill=darkNavy!10, line width=1.2pt]
    \tikzstyle{loaderbox} = [rectangle, rounded corners, minimum width=9cm, minimum height=1cm, text centered, draw=spatialBlue, fill=spatialBlue!15, line width=1.2pt]
    \tikzstyle{runtimebox} = [rectangle, rounded corners, minimum width=9cm, minimum height=1cm, text centered, draw=openxrTeal, fill=openxrTeal!15, line width=1.2pt]
    \tikzstyle{hwbox} = [rectangle, rounded corners, minimum width=9cm, minimum height=1cm, text centered, draw=amberGlow, fill=amberGlow!15, line width=1.2pt]
    \tikzstyle{arrow} = [thick, ->, >=stealth]

    % Nodes
    \node (app) [appbox] {\textbf{STEM Application Layer} (Unity, Unreal, WebXR, Python Engine)};
    \node (api) [below of=app, node distance=1.4cm] {\textbf{OpenXR API} (Function Pointers \& Khronos Registry)};
    \node (loader) [loaderbox, below of=api, node distance=1.2cm] {\textbf{OpenXR Loader} (Validation Layers, API Interception)};
    \node (runtime) [runtimebox, below of=loader, node distance=1.5cm] {\textbf{OpenXR Runtime} (Monado, SteamVR, Meta XR, Windows MR)};
    \node (hw) [hwbox, below of=runtime, node distance=1.5cm] {\textbf{Physical XR Devices} (Sensors, Displays, Haptics, Tracking)};

    % Connections
    \draw [arrow, line width=1.2pt, darkNavy] (app) -- (loader);
    \draw [arrow, line width=1.2pt, spatialBlue] (loader) -- (runtime);
    \draw [arrow, line width=1.2pt, openxrTeal] (runtime) -- (hw);
\end{tikzpicture}
\caption{แผนผังสถาปัตยกรรม OpenXR แสดงลำดับชั้นการเชื่อมต่อระหว่างแอปพลิเคชัน ตัวโหลด รันไทม์ และอุปกรณ์ฮาร์ดแวร์}
\label{fig:openxr_architecture}
\end{figure}

\begin{stemconceptbox}{การแยกส่วนต่อประสานกับการปฏิบัติงานจริง (API vs. Runtime Decoupling)}
หัวใจสำคัญของ OpenXR คือการแยกส่วนคำสั่งมาตรฐาน (OpenXR API) ออกจากการทำงานจริงของฮาร์ดแวร์ (Runtime Implementation) ทำให้โค้ดของห้องทดลองวิทยาศาสตร์เสมือนจริงไม่ต้องอิงกับชุดคำสั่งเฉพาะของผู้ผลิตรายใดเลย เมื่ออุปกรณ์รุ่นใหม่เปิดตัว ผู้ผลิตเพียงสร้าง OpenXR Runtime ที่รองรับมาตรฐาน ตัวแอปพลิเคชันสะเต็มเดิมก็สามารถทำงานบนอุปกรณ์ใหม่ได้ทันที
\end{stemconceptbox}

\section{วงรอบชีวิตของระบบและเฟรมลูป}

การทำงานของแอปพลิเคชัน OpenXR ดำเนินไปตามลำดับขั้นตอนที่มีความเคร่งครัด เพื่อให้มั่นใจว่าการจัดสรรทรัพยากรฮาร์ดแวร์ การคำนวณตำแหน่งเซนเซอร์ และการเรนเดอร์ภาพมีความสอดคล้องกับเวลาจริง (Real-Time Synchronization) ดังแสดงตามขั้นตอนต่อไปนี้

\begin{enumerate}
    \item \textbf{การสร้างอินสแตนซ์ (Instance Creation):} การเรียกฟังก์ชัน \texttt{xrCreateInstance} เพื่อเตรียมความพร้อมของ Loader และตรวจสอบส่วนขยายที่ต้องการ
    \item \textbf{การระบุระบบฮาร์ดแวร์ (System Retrieval):} การเรียกฟังก์ชัน \texttt{xrGetSystem} เพื่อเลือกระบบฮาร์ดแวร์ที่พร้อมทำงาน เช่น รูปแบบ HMD หรือ Handheld
    \item \textbf{การสร้างเซสชัน (Session Lifecycle):} การเชื่อมต่อระบบจำลองภาพเข้ากับหน้าต่างแสดงผลผ่าน \texttt{xrCreateSession} ซึ่งเซสชันจะมีสถานะต่าง ๆ เช่น \texttt{IDLE}, \texttt{READY}, \texttt{SYNCHRONIZED}, \texttt{VISIBLE}, และ \texttt{FOCUSED}
    \item \textbf{วงรอบการเรนเดอร์เฟรม (Frame Loop):} การประมวลผลต่อเนื่องในแต่ละเฟรม ซึ่งประกอบด้วยการรอสัญญาณเฟรม (\texttt{xrWaitFrame}) การเริ่มต้นเฟรม (\texttt{xrBeginFrame}) และการสิ้นสุดเฟรม (\texttt{xrEndFrame})
\end{enumerate}

\begin{workedexamplebox}{การคำนวณอัตราการเรนเดอร์และงบประมาณเวลาต่อเฟรม}
ในการสร้างห้องทดลองเสมือนจริงที่ให้อรรถรสสมจริงและไม่ก่อให้เกิดอาการเมารถจากโลกเสมือน (Cyber Sickness) ระบบต้องเรนเดอร์ภาพด้วยความถี่ไม่น้อยกว่า $f = 90\text{ Hz}$ จงคำนวณเวลางบประมาณสูงสุดต่อหนึ่งเฟรม (Frame Budget) สำหรับการคำนวณฟิสิกส์และการเรนเดอร์ภาพ

\textbf{วิธีทำ:}
งบประมาณเวลาทั้งหมดต่อหนึ่งเฟรมคำนวณได้จากส่วนกลับของความถี่ในการเรนเดอร์
\begin{equation}
\Delta t = \frac{1}{f} = \frac{1}{90\text{ s}^{-1}} \approx 11.11\text{ ms}
\end{equation}
หากแบ่งเวลาออกเป็นการคำนวณทางฟิสิกส์ 30\% การอนุมานท่าทางมือ 20\% และการเรนเดอร์ภาพ 50\% จะได้งบประมาณเวลาในแต่ละส่วน ดังนี้
\begin{align}
\Delta t_{\text{physics}} &= 0.30 \times 11.11\text{ ms} = 3.33\text{ ms} \\
\Delta t_{\text{tracking}} &= 0.20 \times 11.11\text{ ms} = 2.22\text{ ms} \\
\Delta t_{\text{render}} &= 0.50 \times 11.11\text{ ms} = 5.56\text{ ms}
\end{align}
การบริหารเวลานี้มีความสำคัญยิ่งในการเขียนเอนจินจำลองวิทยาศาสตร์ เพื่อไม่ให้เกิดปัญหากระตุกหรือเฟรมตก
\end{workedexamplebox}

\section{การบูรณาการในสะเต็มศึกษาสมัยใหม่}

ในการศึกษาฟิสิกส์และวิทยาศาสตร์กายภาพ มาตรฐาน OpenXR ช่วยให้นักเรียนสามารถมองเห็นเส้นแรงแม่เหล็กที่ตามองไม่เห็น สามารถใช้มือหยิบจับอนุภาคในระดับอะตอม หรือสังเกตวิวัฒนาการของระบบดาวคู่ได้อย่างเป็นรูปธรรม สอดคล้องกับงานวิจัยของ ทัศนา และคณะ ในการใช้แบบจำลองทางคณิตศาสตร์และการคำนวณเชิงตัวเลขเพื่ออธิบายปรากฏการณ์ทางกายภาพที่มีความซับซ้อนสูง ซึ่งการจำลองแบบเสมือนจริงช่วยยกระดับการรับรู้ของผู้เรียนได้อย่างมีนัยสำคัญ

\begin{aipromptbox}[การสร้างสเกเลตันโค้ด OpenXR Minimal Session ด้วย AI]
\textbf{บริบทและเป้าหมาย:} สร้างโค้ดเริ่มต้นเชื่อมต่อ OpenXR Loader และเริ่ม Session สำหรับแอปพลิเคชันสะเต็มศึกษา\\
\textbf{CLEAR Prompting Template (คัดลอกนำไปใช้งานได้ทันที):}
\begin{lstlisting}[style=openxrcode]
Act as an expert Khronos OpenXR Systems Engineer.
Goal: Write a production-grade Python script using pyopenxr for a STEM Virtual Lab that:
1. Initializes an XrInstance with application name "STEM_Spatial_Lab" version 1.0.
2. Checks available extensions and enables XR_KHR_opengl_enable and XR_EXT_hand_tracking.
3. Retrieves the system ID for XR_FORM_FACTOR_HEAD_MOUNTED_DISPLAY.
4. Creates an XrSession and manages the state machine (READY, SYNCHRONIZED, FOCUSED).
5. Implements an explicit 90 FPS frame loop with error handling and clean resource destruction.
Output: Complete executable Python code with informative English/Thai comments.
\end{lstlisting}
\end{aipromptbox}

\begin{codebox}[การตรวจสอบส่วนขยายที่รองรับใน OpenXR Runtime ด้วย Python]
import openxr as xr

def check_openxr_runtime():
    # ตรวจสอบส่วนขยายทั้งหมดที่รันไทม์รองรับ
    extensions = xr.enumerate_instance_extension_properties()
    print(f"ตรวจพบส่วนขยายทั้งหมด {len(extensions)} รายการ:")
    for ext in extensions:
        print(f" - {ext.extension_name} (v{ext.extension_version})")

    # ตรวจสอบส่วนขยายจำเป็นสำหรับสะเต็มศึกษา
    required = ["XR_KHR_opengl_enable", "XR_EXT_hand_tracking"]
    for req in required:
        supported = any(ext.extension_name == req for ext in extensions)
        status = "พร้อมใช้งาน" if supported else "ไม่รองรับ"
        print(f"ส่วนขยาย {req}: {status}")

if __name__ == "__main__":
    check_openxr_runtime()
\end{codebox}

\section*{สรุปท้ายบทที่ 1}
\addcontentsline{toc}{section}{สรุปท้ายบทที่ 1}


มาตรฐาน OpenXR ได้สร้างรากฐานสำคัญสำหรับการพัฒนาเทคโนโลยีความเป็นจริงขยาย ด้วยการสร้างสถาปัตยกรรมที่เป็นเอกภาพระหว่างแอปพลิเคชันและฮาร์ดแวร์ การทำความเข้าใจโครงสร้าง วงรอบชีวิตของระบบ และข้อจำกัดด้านเวลาในการเรนเดอร์ เป็นขั้นตอนแรกที่จำเป็นอย่างยิ่งในการออกแบบและพัฒนาห้องทดลองวิทยาศาสตร์และสะเต็มศึกษาเสมือนจริงที่มีประสิทธิภาพสูงในบทเรียนถัดไป

\section*{แบบฝึกหัดท้ายบทที่ 1}
\addcontentsline{toc}{section}{แบบฝึกหัดท้ายบทที่ 1}

\begin{enumerate}
    \item จงอธิบายความแตกต่างระหว่าง OpenXR API และ OpenXR Runtime พร้อมยกตัวอย่างชื่อรันไทม์ที่ใช้กันอย่างแพร่หลาย
    \item เพราะเหตุใดจึงต้องมีการจำแนกสถานะของ OpenXR Session ออกเป็นหลายระดับ เช่น \texttt{VISIBLE} และ \texttt{FOCUSED}
    \item หากต้องการพัฒนาห้องปฏิบัติการเสมือนจริงบนอุปกรณ์ที่ทำงานด้วยอัตราการรีเฟรช $120\text{ Hz}$ จงคำนวณงบประมาณเวลาต่อเฟรม และเสนอแนวทางการจัดสรรเวลาสำหรับการคำนวณฟิสิกส์และการตรวจจับการเคลื่อนไหว
\end{enumerate}
